% paper/sections/11d_time_accuracy.tex
%
% §11.3  NS 方程式の総合精度
% ═══════════════════════════════════════════════════════════════════════════════
\subsection{NS 方程式の総合精度}
\label{sec:ns_time_accuracy}
% ═══════════════════════════════════════════════════════════════════════════════

コンポーネント単体（第~\ref{sec:numerical_integrity}章）ではなく，
Predictor--PPE--Corrector パイプライン全体としての時間・空間収束次数を検証する．

% ─────────────────────────────────────────────────────────────────────────────
\subsubsection{Taylor--Green 渦による時間収束次数}
\label{sec:tgv_temporal}
% ─────────────────────────────────────────────────────────────────────────────

\textbf{設定：}
§\ref{sec:energy_conservation} と同一の Taylor--Green 渦を
$N = 64$ に固定（空間誤差を十分に小さくする）し，
$\Delta t$ を系統的に変化させて $T = 0.5$ における速度場の $L^\infty$ 誤差を計測する．

\begin{table}[htbp]
\centering
\caption{NS 時間精度：Taylor--Green 渦（$N=64$，$\mathit{Re}=100$，AB2+IPC）}
\label{tab:tgv_temporal}
\begin{tabular}{rrrc}
\toprule
ステップ数 & $\Delta t$ & $L^\infty(\bu)$ & 次数 \\
\midrule
  5  & $0.1000$ & $3.35 \times 10^{-9}$  & --- \\
 10  & $0.0500$ & $8.37 \times 10^{-10}$ & $2.00$ \\
 20  & $0.0250$ & $2.10 \times 10^{-10}$ & $1.99$ \\
 40  & $0.0125$ & $5.33 \times 10^{-11}$ & $1.98$ \\
 80  & $0.0063$ & $1.42 \times 10^{-11}$ & $1.91$ \\
160  & $0.0031$ & $4.34 \times 10^{-12}$ & $1.71$ \\
\bottomrule
\end{tabular}
\end{table}

\textbf{結果：}
$\Delta t \geq 0.0125$ の領域で精確な $\Ord{\Delta t^2}$ 収束が確認された（次数 $1.98$--$2.00$）．
これは AB2 時間外挿（$\Ord{\Delta t^2}$）+ IPC 分割誤差（$\Ord{\Delta t^2}$，van Kan 1986）
の理論的期待と完全に一致する．
$\Delta t \leq 0.006$ では丸め誤差のフロア（$\sim 10^{-12}$）に接近し，見かけ上の次数が低下する．

% ─────────────────────────────────────────────────────────────────────────────
\subsubsection{Kovasznay 流による空間精度の検証}
\label{sec:kovasznay}
% ─────────────────────────────────────────────────────────────────────────────

\textbf{設定：}
Kovasznay (1948) の 2 次元定常 NS 解析解（$\mathit{Re} = 40$）を用いる：
\begin{align}
  u(x,y) &= 1 - e^{\lambda x}\cos(2\pi y), &
  v(x,y) &= \frac{\lambda}{2\pi}\,e^{\lambda x}\sin(2\pi y), \\
  p(x,y) &= -\tfrac{1}{2}\,e^{2\lambda x},
\end{align}
ただし $\lambda = \mathit{Re}/2 - \sqrt{\mathit{Re}^2/4 + 4\pi^2} \approx -6.42$．
計算領域 $[0, 1] \times [-0.5, 0.5]$，ディリクレ境界条件．

CCD で離散化した NS 残差 $R = -\bu \cdot\bnabla\bu + \nu\nabla^2\bu - \bnabla p$ を
解析解上で評価し，格子精緻化に伴う残差の減少から空間精度を測定する
（MMS 残差法）．

\begin{table}[htbp]
\centering
\caption{Kovasznay 流の CCD 離散化 NS 残差（$\mathit{Re}=40$，ディリクレ BC）}
\label{tab:kovasznay}
\begin{tabular}{rrcrc}
\toprule
$N$ & NS 残差 $\|R\|_\infty$ & 次数 & $\|\bnabla\cdot\bu\|_\infty$ & 次数 \\
\midrule
 16 & $2.60 \times 10^{-4}$ & --- & $6.60 \times 10^{-5}$ & --- \\
 32 & $2.28 \times 10^{-5}$ & $3.51$ & $1.24 \times 10^{-6}$ & $5.74$ \\
 64 & $1.55 \times 10^{-6}$ & $3.88$ & $2.05 \times 10^{-8}$ & $5.91$ \\
128 & $9.93 \times 10^{-8}$ & $3.96$ & $3.27 \times 10^{-10}$ & $5.97$ \\
\bottomrule
\end{tabular}
\end{table}

\textbf{結果：}
CCD 離散化連続方程式の残差は $\Ord{h^{5.97}}$ で収束し，CCD 内部精度 $\Ord{h^6}$ に整合する．
NS 運動量残差は $\Ord{h^{3.96}}$ に収束する．これは対流項 $\bu \cdot \bnabla\bu$
が非線形積であるため，ディリクレ境界近傍で CCD 境界スキーム（$\Ord{h^5}$）の
誤差が非線形に増幅されることによる．
内部節点のみを評価しても境界スキームの影響は数節点にわたり伝播するため，
$\Ord{h^4}$ 程度の実効精度となる．
$N=128$ で NS 残差は $10^{-8}$ に達しており，
CCD ベースの NS ソルバが十分な空間精度を有することが確認された．
